\section{Motivation and Contributions}
\label{thesis:intro-motivation}

Recent trends in cloud computing shifted the paradigm from running workloads on
local premises to shared infrastructures where multiple tenants share the same
physical resources for higher scalability and lower costs. In addition, with the
\acl{IoT} and edge computing, some trusted services need to be offloaded to
devices that may be in control of a malicious user, like a car or a smartphone.
Here, \acp{TEE} can play an important role to isolate sensitive workloads even
on compromised nodes, and indeed this technology is becoming more and more
ubiquitous in many different fields and use cases. Compared to other approaches
such as \ac{FHE}~\cite{acar2018fhe} and \ac{MPC}~\cite{lindell2020mpc}, which
are active areas of research but do not still have many real-world applications,
\acp{TEE} are a more accessible solution that can be integrated into existing
workloads with relatively low effort and performance overhead.

However, there are some operational aspects that should be considered when using
\acp{TEE}, such as how to securely provision and attest workloads, how to
protect and authenticate communication between different services, and how to
deal with peripherals such as sensors and actuators. Besides, there are
differences between running an application on a local, closed system and running
it within a \ac{TEE} on a third-party infrastructure such as a public cloud. For
example, applications that rely on system resources, such as a clock, may become
vulnerable to attacks if running on a shared system, because a privileged
adversary can control these resources and may thus influence the behavior of the
application. Therefore, moving workloads inside \acp{TEE} requires careful study
and, possibly, some application-level changes. This dissertation delves deeper
into these challenges, presenting solutions to securely integrating \acp{TEE}
into new use cases and acknowledging limitations of this technology.

%Note that we focus on \emph{operational} aspects of \acp{TEE}, i.e., how to use
%them in the ``correct'' way as they are designed to, and how to simplify
%operations such as provisioning and remote attestation.

Below, we describe the contributions of this Ph.D. trajectory. All the code
developed during our work has been open-sourced for further research and
reproducibility of results. This involves proof of concepts, frameworks,
prototypes and tools. For each contribution summarized below, we link to the
public codebase (if any).